%============================================================
% Chapter 16: GCD, Divisibility & Number-Theoretic Arrays
% STATUS: PLANNED  (generated by tools/outline.py; replace this file when the chapter is written)
%============================================================
\chapter{GCD, Divisibility \& Number-Theoretic Arrays}
\label{ch:number_arrays}

\begin{tcolorbox}[enhanced, colback=white, colframe=secondaryblue, boxrule=1pt, arc=2mm]
\textbf{\large\textcolor{primaryblue}{Arrays of integers carry arithmetic structure: prefix GCDs, remainders, divisor counts and harmonic-sum loops over values.}}
\end{tcolorbox}

\chapterfacts{1500}{2300}{6}{12}{30}

\begin{plannedbox}[Planned Contents --- this chapter is not written yet]
\textbf{Part A --- Lessons}
\begin{enumerate}[leftmargin=*, label=\textbf{A\arabic*.}, itemsep=1pt]
    \item Prefix and suffix GCD; the GCD of a subarray changes $\bigO{\log V}$ times
    \item Remainders: counting pairs and subarrays modulo $k$
    \item Loops over multiples: $\sum V/d = \bigO{V \log V}$ and counting by divisors
    \item Sieve-based preprocessing of the value domain
    \item Inclusion--exclusion and M\"obius over the values of an array
    \item The pigeonhole principle on prefix sums and remainders: when a large $n$ forces the answer
\end{enumerate}
\textbf{Part B --- Patterns} (with candidate problems, checked when the chapter is written)
\begin{enumerate}[leftmargin=*, label=\textbf{P\arabic*.}, itemsep=1pt]
    \item \textbf{Prefix and suffix GCD}
    \item \textbf{Distinct GCDs of subarrays}
    \item \textbf{Remainder counting}
    \item \textbf{Harmonic loops over values}
    \item \textbf{Counting coprime pairs}
    \item \textbf{Pigeonhole: large inputs force the answer}
\end{enumerate}
\end{plannedbox}
